但资源量回升并非等同于生态系统向良性方向恢复。部分水域鱼类增长过快，对水生植物和浮游动物形成过度摄食，继而造成水体浑浊和底栖植被减少，部分水域甚至形成大面积"水下荒漠"。与此同时，长江江豚数量已回升至约1249头，2024年重新发现长江鲟自然繁殖产卵场；而鳄雀鳝、巴西龟等外来物种已在鄱阳湖等流域支流存在稳定种群。十年禁渔的资源恢复成效确实存在，系统却未必更加健康，营养级失衡、连通性不足和入侵扩散依然突出。因此有必要分别评价"资源恢复"与"系统健康"。
\cite{report2022,report2023}。
